\documentclass[../../main.tex]{subfiles}
\begin{document}

{\bf [解答]}
\paragraph{(1)}

与えられた関数
\begin{align}
  f(x) = \sin(\sin x) \label{eq:1}
\end{align}
について，
\begin{align}
  f(x+2\pi) = \sin(\sin(x+2\pi)) = \sin(\sin x) = f(x)
\end{align}
から$f(x)$は周期関数である. $2\pi$は周期の一つである．

そこで以下周期が$2\pi$であることを示す.
$0 < c < 2\pi, \ f(x+c) = f(x)$をみたす$c$があるとすると
\begin{align}
  \sin(x+c) = \sin x + 2n\pi, \ -\sin x + (2n+1)\pi \quad (n \in \mathbb{Z})
\end{align}
と表せるが, $-1 \le \sin(x+c), \sin x \le 1, \ \pi \simeq 3.14$から, これらをみたす$n$は$0$のみであるから
\begin{align}
  \sin(x+c) = \pm \sin x
\end{align}
となる．これを任意の$x$でみたす$c \ (0 < c < 2\pi)$は存在せず, 矛盾する．
よって周期は$2\pi$である.

\paragraph{(2)}

与えられた関数
\begin{align}
  f(x) = \cos(\sin x)
\end{align}
について
\begin{align}
  f(x+2\pi) = \cos(\sin (x+2\pi)) = \cos(\sin x) = f(x)
\end{align}
より$f(x)$は周期関数で，$2\pi$は周期の一つである．

そこで以下周期が$2\pi$であることを示す. $0 < c < 2\pi, \ f(x+c) = f(x)$をみたす$c$があるとすると
\begin{align}
  & \cos(\sin(x+c)) = \cos(\sin x) \\
  \therefore
  & \cos(x+c) = \cos x + 2n\pi, \ -\cos x + 2n\pi \quad (n \in \mathbb{Z})
\end{align}
である．$-1 \le \cos(x+c), \cos x \le 1$から, $n=0$が必要で
\begin{align}
  \cos(x+c) = \pm \cos x
\end{align}
となる．これを任意の$x$でみたす$c$は存在せず, 矛盾する．
よって周期は$2\pi$である.

\paragraph{(3)}

与えられた関数
\begin{align}
  f(x) = \sin (x^3)
\end{align}
について，$f(x+c) = f(x)$をみたす$c \ (0 < c)$が存在するとする.
まず$x=0$での成立が必要で,
\begin{align}
  \sin c^3 = \sin 0 = 0 \\
  \therefore
  c = \sqrt[3]{n\pi} \quad (n \in \mathbb{N})
\end{align}
とおける.

さらに $0 < x < \sqrt[3]{\pi}$の時, $h(x) \neq 0$だから$c < x < c + \sqrt[3]{\pi}$でも$h(x) \neq 0$であることが必要だが, $x = \sqrt[3]{(n+1)\pi}$はこの区間内にあるが$h(\sqrt[3]{(n+1)\pi}) = 0$だから矛盾する.
よって$f(x)$は周期関数ではない.

{\bf [解説]}

(1), (2)は $-\dfrac{\pi}{2} < \sin x < \dfrac{\pi}{2}$を使うと早い．

\end{document}
